\documentclass[12pt,a4paper]{ctexart}

\usepackage[margin=2.5cm]{geometry}
\usepackage{hyperref}

\hypersetup{
  hidelinks
}

\setlength{\parskip}{0.45em}
\setlength{\parindent}{2em}

\begin{document}

\begin{center}
{\LARGE\heiti 测试计划}
\end{center}

\section{黑盒测试计划}

\subsection{黑盒测试目标}

黑盒测试重点检查主要功能、角色权限和核心业务流程是否符合需求。具体目标如下：

\begin{enumerate}
  \item 验证系统主要功能是否与需求描述一致。
  \item 验证教师、学生两类角色的功能权限是否正确。
  \item 验证课程平台核心业务流程能否端到端顺利完成。
  \item 基于等价类、边界值、决策表和场景法设计测试用例。
  \item 使用自动化接口测试与少量人工页面测试相结合的方式完成验证。
\end{enumerate}

\subsection{黑盒测试范围与重点}

被测对象为当前仓库中的 \texttt{RAG-Knowledge-Assistant} 项目。该项目包含前端与后端两部分，功能覆盖用户管理、课程管理、课程资料管理、练习题管理和课程社区管理，满足课程作业对项目规模和功能点数量的要求。

本次黑盒测试覆盖以下功能点：

\begin{enumerate}
  \item 用户注册
  \item 用户登录
  \item 个人信息查看
  \item 个人信息修改
  \item 教师创建课程
  \item 学生加入课程
  \item 学生退出课程
  \item 课程资料上传
  \item 教师解散课程
  \item 课程资料查看
  \item 课程资料下载
  \item 课程资料 PDF 预览
  \item 课程资料删除
  \item 练习题上传
  \item 练习题查看
  \item 练习题作答
  \item 作答结果反馈
  \item 课程社区发帖
  \item 课程社区评论
  \item 课程社区回复
  \item 核心流程黑盒验证
\end{enumerate}

重点检查以下内容：

\begin{enumerate}
  \item 不同角色的权限控制是否正确。
  \item 各业务操作后的页面表现与数据状态是否一致。
  \item 文件相关功能是否正确处理格式、大小限制和异常情况。
  \item 练习题与课程社区功能是否能在课程上下文中正确关联。
  \item 核心业务流程中的前后步骤是否能够正确衔接。
\end{enumerate}

核心流程定义为：

\begin{quote}
教师登录 $\rightarrow$ 创建课程 $\rightarrow$ 上传资料 $\rightarrow$ 上传练习题 $\rightarrow$ 学生登录 $\rightarrow$ 加入课程 $\rightarrow$ 查看资料 $\rightarrow$ 完成练习 $\rightarrow$ 获得反馈 $\rightarrow$ 参与课程社区发帖 / 评论 / 回复
\end{quote}

\subsection{黑盒测试方法}

测试条件和测试用例主要采用以下方法设计：

\subsubsection{等价类划分}

将输入条件与业务状态划分为有效等价类和无效等价类，主要包括：

\begin{enumerate}
  \item 角色：教师、学生、非法角色。
  \item 登录信息：正确账号密码、错误密码、未注册账号、空输入。
  \item 文件类型：PDF、非 PDF、ZIP、超限文件。
  \item 课程状态：未加入、已加入、已解散、非本人课程。
\end{enumerate}

\subsubsection{边界值分析}

重点覆盖以下边界条件：

\begin{enumerate}
  \item 分页页码边界。
  \item 文件大小上限边界。
  \item 标题、描述、评论内容等空输入边界。
  \item ZIP 压缩包单层目录与多层目录边界。
  \item 练习题集题目数量边界。
\end{enumerate}

\subsubsection{决策表法}

主要用于角色权限和条件组合较明显的场景，例如：

\begin{enumerate}
  \item 注册是否允许提交。
  \item 登录是否成功。
  \item 创建课程是否允许执行。
  \item 加入课程、退出课程、解散课程是否满足角色与状态条件。
  \item PDF 预览与评论回复是否满足前置条件。
\end{enumerate}

\subsubsection{场景法}

用于验证跨模块核心流程，检查真实教学场景下业务操作能否连续完成，以及模块之间的数据关联与状态流转是否正确。

\subsection{黑盒测试环境、工具与数据}

\subsubsection{测试环境}

\begin{enumerate}
  \item 操作系统：Windows 11 开发环境。
  \item 前端：React 19。
  \item 后端：FastAPI。
  \item 浏览器：Chrome 与 Edge。
  \item 前端访问地址：\url{http://localhost:3000}。
  \item 后端接口地址：\url{http://127.0.0.1:8000}。
\end{enumerate}

\subsubsection{测试工具}

\begin{enumerate}
  \item \texttt{Codex}：用于辅助梳理需求、编写测试计划、设计测试用例、生成自动化测试代码和整理测试结果。
  \item \texttt{pytest + FastAPI TestClient}：用于执行自动化黑盒接口测试。
  \item 浏览器人工操作：用于补充自动化测试难以完全覆盖的页面交互、PDF 预览、下载行为和页面级流程验证。
\end{enumerate}

\subsubsection{测试数据}

\begin{enumerate}
  \item 教师账号 1 个。
  \item 学生账号 1 个。
  \item 测试课程若干。
  \item 测试资料若干，包含 PDF、非 PDF、ZIP 等文件。
  \item 测试练习题集若干，覆盖单选、多选、填空三类题型。
\end{enumerate}

\subsection{黑盒测试执行安排}

黑盒测试采用“自动化接口测试为主、少量人工页面测试为辅”的方式执行。自动化测试承担主体覆盖任务，人工测试主要补充前端页面交互、文件下载与 PDF 预览等浏览器行为验证。

执行顺序如下：

\begin{enumerate}
  \item 环境与测试数据准备。
  \item 用户管理模块测试。
  \item 课程管理模块测试。
  \item 课程资料管理模块测试。
  \item 练习题管理模块测试。
  \item 课程社区管理模块测试。
  \item 核心流程端到端测试。
  \item 回归验证与结果整理。
\end{enumerate}

测试优先级划分如下：

\begin{enumerate}
  \item \texttt{P0}：用户登录、教师创建课程、学生加入课程、资料上传、练习题查看与作答、核心流程贯通。
  \item \texttt{P1}：用户注册、个人信息修改、资料下载、PDF 预览、社区发帖评论回复。
  \item \texttt{P2}：课程搜索分页、资料列表分页、帖子搜索分页等辅助功能。
\end{enumerate}

\section{白盒测试计划}

\end{document}
